\documentstyle[%


righttag%


,11pt%


,amssymb%


,russian%               remove in Eng.


,izvru%                 Set izven% in Eng.


]{amsart}





%  separating commas!





\newcommand{\tts}[1]{}





%\hoffset=-15mm%                  For Eng.


\setcounter{page}{15}


\god{2000}


\nomer{12~(463)} %                        Remove in Eng.


%\nomer{Vol.\,44, No.\,12}%               Set in Eng.


%\str{\pageref{begin}--\pageref{end}}%  Set in Eng.


\udk{519.712}





\author{\it И.П.\,Вознюк}


% NB:   ^^ remove in Eng.


\title{Аппроксимационный алгоритм для задачи размещения на максимум


с ограниченными объемами производств и поставок}


\thanks{Работа выполнена при финансовой поддержке 


федеральной целевой программы


``Интеграция'' (код проекта 274).}





\date{}


%\pagestyle{myheadings}%         Set in Eng.


\pagestyle{plain} %              Remove in Eng.


\begin{document}


%\thispagestyle{plain}%          Set in Eng.


\maketitle


\label{begin}





\section{Введение}





Рассматривается следующая задача размещения на максимум:


дано множество $I$ --- возможных пунктов размещения предприятий,


$J$ --- множество пунктов потребления.


Известны штрафы за размещения предприятий в пунктах $I$


и объемы спроса в пунктах потребления $J$.


Известна также прибыль, получаемая при поставке 


единицы продукта от производителя к потребителю.


Заданы ограничения на объемы производств


и ограничения на объемы поставок продукта от производителей 


к потребителям. 


Требуется на множестве $S \subseteq I$ разместить предприятия


так, чтобы максимизировать общую прибыль, получаемую при удовлетворении 


потребностей в продукте по всем пунктам потребления с учетом штрафов 


за размещения предприятий.





Обобщением ограничений на объемы поставок являются ограничения на


пропускные способности коммуникаций для сетевой задачи размещения на 


минимум \cite{V99}. В указанной статье предложен 


псевдополиномиальный алгоритм для решения задачи


на древесной сети с ограничениями на пропускные способности


коммуникаций.





Исходная задача $NP$-трудная, т.\,к. к задаче размещения 


без дополнительных ограничений сводится $NP$-трудная задача о вершинном 


покрытии (\cite{GD82}, с.\,74).


Поэтому представляется интересным вопрос о нахождении приближенных 


решений этой задачи.





Был построен ряд аппроксимационных алгоритмов 


для задач размещения на минимум.


Для общего случая задачи 


известен полиномиальный жадный алгоритм, отыскивающий допустимое 


решение задачи, отличающееся от точного не более чем в 


$\ln n$ раз, где $n$ --- число пунктов потребления \cite{H82}.


В статье \cite{STA97} для случая, когда


транспортные затраты удовлетворяют неравенству треугольника,


предложен алгоритм,


который находит допустимое решение, отличающееся от точного не более


чем в $3{,}16$ раз.


Вслед за этой работой появились результаты, улучшающие эту оценку для 


рассматриваемой задачи. В \cite{GK98} и 


\cite{C98} предложены аппроксимационные 


алгоритмы, находящие приближенные решения задачи с точностью $2{,}41$


и $1{,}74$ соответственно. 





Для задачи размещения на минимум с ограниченными объемами производств


известны следующие результаты. 


В \cite{STA97} для задачи с одинаковыми по величине 


объемами производств предложен аппроксимационный алгоритм,


который находит допустимое решение, отличающееся от точного не более


чем в $7$ раз.


При этом допускается увеличить заданные объемы производств в $3{,}5$ раза.


В \cite{KPR98} и


\cite{CW99} предлагаются алгоритмы, которые находят приближенные 


решения этой задачи


с точностью $(8+ \epsilon)$ и $6(1+\epsilon)$ соответственно


для любого $\epsilon >0$ без увеличения мощности предприятий. 





Для задачи размещения на максимум в \cite{CFN77}


построен жадный алгоритм,


который находит допустимое решение со следующей точностью


для относительной оценки со сдвигом:


$(Z^*-Z_G)/(Z^*-Z_R) < 1/e\approx 0{,}368$,


где $Z^*$ --- значение оптимального решения,


$Z_G$ --- величина допустимого решения, полученного предлагаемым алгоритмом,


$Z_R$ --- гарантированная нижняя оценка решения,


$e$ --- основание натурального логарифма.


В \cite{AS99} улучшен данный результат и предлагается алгоритм с 


оценкой точности 


$(Z^*-Z_G)/(Z^*-Z_R) \le 3- 2\sqrt{2} 


\approx 0{,}172$.





В данной работе 


показывается, что относительная оценка со сдвигом, равная $1/e$,


гарантируется жадным алгоритмом 


для задачи на максимум с ограничениями на 


объемы производств и поставок.





\section{Постановка задачи}





Пусть $I = \{1, \dotsc, m \}$ и $J = \{1, \dotsc, n \}$.


Пусть заданы штрафы $f_{i} \geq 0, i \in I$, за размещение предприятия 


и объемы $b_{j} \ge 0, j \in J$, потребления продукта.


Указана прибыль $g_{ij}$, получаемая при


доставке единицы продукта из $i$ в $j$, 


заданы ограничения на объемы производств $d_i$ и


ограничения на максимальный объем поставки продукта


$a_{ij}$, который может 


поставляться из $i$ в $j$, $i \in I$, $j \in J$. 


Требуется указать те предприятия из $I$, 


при вводе которых в систему,


во-первых, полностью бы удовлетворялись спросы 


всех пунктов потребления при соблюдении ограничений на объемы 


производств и поставок и, во-вторых, 


прибыль, получаемая при обслуживании всех пунктов потребления


с учетом штрафов за размещения предприятий, была бы максимальной.





Пусть переменная $x_{ij}$ равна количеству единиц продукта, поставляемого из


$i$ в $j$, $i \in I$, $j \in J$.


Задача размещения на максимум с ограниченными объемами производств и 


поставок формулируется в следующем виде: 


\begin{gather}


\label{l1}


\max \limits_{S \subseteq I} Z(S), \\


%


\label{l10}


Z(S)= \max_{(x_{ij})} \sum_{i \in S} \sum_{j \in J} g_{ij}x_{ij} -


\sum_{i \in S} f_i, \\


%


\label{l2}


\sum _{i \in S} x_{ij}=b_j,\ \; j \in J, \\


%


\label{l201}


\sum _{j \in J} x_{ij} \le d_i,\ \; i \in S, \\


%


\label{l3}


0\le x_{ij}\le a_{ij},\ \; x_{ij}\text{\; целое, \;}\ \; i\in S,\ \; i\in J.


\end{gather}





\section{Субмодулярные функции и жадный алгоритм}





В этом разделе, следуя \cite{NWF78},


дадим определение и напомним некоторые свойства 


субмодулярных функций. Далее опишем жадный 


алгоритм, отыскивающий приближенное решение задачи, целевая функция которой 


обладает свойством субмодулярности.


И, наконец, представим гарантированную 


оценку точности решения, получаемого этим алгоритмом.





\begin{defn}


Пусть $I$ --- конечное множество. Вещественная функция $Z$ на множестве


подмножеств $I$ называется {\it субмодулярной}, если


выполняется условие


$$


Z(A)+Z(B) \ge Z(A \cup B) +Z(A \cap B)\ \; \forall A,B \subseteq I.


$$


\end{defn}





Справедливы три утверждения, полученные в \cite{NWF78}.





\begin{lem} 


\label{c2}


Пусть $d_i$, $i \in I$, --- произвольные вещественные числа. Тогда линейная 


функция $z(S)=\sum\limits_{i \in S}d_i$, $S \subseteq I$,


является субмодулярной.


\end{lem}





\begin{lem} 


\label{c3}


Положительная линейная


комбинация субмодулярных функций является субмодулярной функцией.


\end{lem}





\begin{lem} 


\label{c4}


Пусть дана субмодулярная неубывающая функция $v$ на множестве подмножеств 


$D$ и семейство подмножеств $\{ D_i \}$, $i \in I$, множества $D$.


Тогда функция 


$z(S) = v( \cup_{i \in S} D_i)$ является субмодулярной функцией на множестве 


подмножеств $I$.


\end{lem}





Рассмотрим задачу


\begin{equation}


\label{c1}


\max_{S \subseteq I} \{Z(S) : |S| \le k,


\ Z(S) \text{ --- субмодулярная функция} \}.


\end{equation}





Обозначим $\rho_i(S)= Z(S \cup \{ i \}) - Z(S)$.


Для отыскания допустимого решения задачи $(\ref{c1})$ в


\cite{NWF78} предложен следующий жадный алгоритм A.





{\it Инициализация}. Пусть $S^0= \emptyset$, $I^0=I$, $t=1$.





{\it Итерация $t$}. Найти $i(t) \in I^{t-1}$ такое, чтобы 


$\rho_{i(t)}(S^{t-1})= \max\limits_{i \in I^{t-1}} \rho_i(S^{t-1})$.


Пусть $\rho_{t-1}= \rho_{i(t)}(S^{t-1})$.





{\it Шаг} 1. Если $\rho_{t-1} \le 0$, то 


алгоритм заканчивает работу с множеством $S^{k^*}$, $k^*=t-1 < k$.


Если $\rho_{t-1} > 0$, то пусть $S^t=S^{t-1} \cup \{ i(t)\}$,


$I^t= I^{t-1} \setminus \{i(t) \}$, и алгоритм продолжaeт работу.





{\it Шаг} 2. Если $t=k$, то алгоритм заканчивает работу с множеством 


$S^{k^*}$, $k^*=k$. Иначе $t=t+1$.





Имеем $Z_G=Z(S^{k^*})$ --- приближенное решение, полученное алгоритмом.


Пусть $Z^*$ --- оптимальное значение задачи $(\ref{c1})$.


Обозначим через $C(\Theta)$, где $\Theta \ge 0$, класс субмодулярных


функций, удовлетворяющих условию 


$\rho_i(S) \ge - \Theta$ $\forall S \subseteq I$, 


$i \in I \setminus S$. Пусть $Z( \emptyset) \ge 0$.





Из работы \cite{NWF78} следует





\begin{thm} 


\label{t1}


Жадный алгоритм $\mathrm A$ в задаче $(\ref{c1})$ при условии, 


что $Z(S) \in C(\Theta)$, приводит к


неравенству


$$


\frac{Z^* -Z_G}{Z^*-Z(\emptyset)+k\Theta} \le (1- 1 / k )^k< 1/e,


$$


где $e$ --- основание натурального логарифма.


\end{thm}





\section{Эквивалентная задача}





Пусть $F= \max\limits_{i \in I}f_i+1$.


Воспользуемся методом штрафных функций для 


построения приближенного решения.





Рассмотрим задачу


\begin{gather}


\label{l21}


\max_{S \subseteq I} H(S), \\


%


\label{l210}


 H(S)= \max_{(x_{ij})}


\sum_{i \in S} \sum_{j \in J} (g_{ij}+F)x_{ij} - 


\sum_{i \in S} f_i- F\sum_{j \in J} b_j, \\


%


\label{l22}


\sum_{i \in S}x_{ij} \le b_j,\ \; j \in J, \\


%


\label{l220}


\sum _{j \in J} x_{ij} \le d_i,\ \; i \in S, \\


%


\label{l23}


0\le x_{ij}\le a_{ij},\ \; x_{ij}\ \;\text{ целое, \;}\ \;i\in S,\ \;j\in J.


\end{gather}





Целевая функция (\ref{l210}) получается после введения штрафной функции


$$


\sum_{j \in J}\bigg( \sum_{i \in I}x_{ij} -b_j\bigg)F


$$ 


в целевую функцию (\ref{l1})


исходной задачи и соответствующих преобразований.


Заметим, что в ограничении $(\ref{l22})$ имеем нестрогое неравенство.





Представленная задача является вспомогательной при отыскании допустимого


решения задачи $(\ref{l1})$--$(\ref{l3})$ с гарантированной оценкой


точности получаемого решения. 


Следующие


вспомогательные утверждения устанавливают связь между задачами


$(\ref{l1})$--$(\ref{l3})$ и


$(\ref{l21})$--$(\ref{l23})$.





\begin{lem}


\label{l4}


Задачи $(\ref{l1})$--$(\ref{l3})$ и 


$(\ref{l21})$--$(\ref{l23})$ эквивалентны.


\end{lem}





\begin{pf}


Введение штрафной функции гарантирует выполнение


равенства в неравенстве $(\ref{l22})$ для всех $i\in I$ в оптимальном решении


задачи $(\ref{l21})$--$(\ref{l23})$.


Допустим, что это не так. 


Пусть $Z^*$ --- значение оптимального решения и пусть для некоторых


$j \in J$ выполняется строгое неравенство в условии $(\ref{l22})$.


Тогда рассмотрим такое допустимое решение со значением 


$Z_A$, что в некоторый 


такой пункт потребления $j$ поставляется еще одна единица продукта 


из некоторого предприятия $i$, которое можно открыть, если $i \notin S$.


Имеем $Z_A > Z^*$, т.\,к.


$Z_A - Z^* \ge g_{ij} -f_i +F > 0$.


Получили решение со значением целевой функции, большим $Z^*$. Противоречие.


Таким образом, оптимальное решение задачи 


$(\ref{l21})$--$(\ref{l23})$ является оптимальным решением задачи


$(\ref{l1})$--$(\ref{l3})$.


Обратное утверждение очевидно. 


\end{pf}





\begin{lem}


\label{l41}


Допустимое решение задачи $(\ref{l21})$--$(\ref{l23})$, полученное 


жадным алгоритмом $\mathrm A$,


является допустимым решением задачи $(\ref{l1})$--$(\ref{l3})$.


\end{lem}





\begin{pf}


Найдя допустимое решение задачи $(\ref{l21})$--$(\ref{l23})$


жадным алгоритмом A, будем иметь равенство в неравенстве $(\ref{l22})$


для всех $j \in J$.


Допустим, что это не так. Пусть $Z_A$ --- значение допустимого решения, 


полученного жадным алгоритмом A, и пусть для некоторых


$j \in J$ выполняется строгое неравенство в условии $(\ref{l22})$.


Тогда найдется такое $i \in I$, что $\rho_i(S) \ge g_{ij} -f_i +F > 0$


для некоторого $j \in J$. Получили противоречие с условием остановки 


алгоритма A. 


Таким образом, алгоритм A получает решение, удовлетворяющее условиям


$(\ref{l2})$--$(\ref{l3})$.


\end{pf}





\section{Оценки качества алгоритма}





Поскольку значение целевой функции представленной задачи


может быть как положительным, так и 


отрицательным, то, следуя \cite{CFN77}, будем


использовать оценку относительного уклонения 


$(Z^*-Z_G)/(Z^*-Z_R)$, где 


$Z^*$ и $Z_G$ --- оптимальное решение и 


допустимое решениe задачи, получаемые предложенным алгоритмом A; 


$Z_R$ --- гарантированная нижняя оценка для решения.


В данной работе


будем полагать


\begin{equation}


Z_R= \sum_{j \in J}b_j \min_{i \in I}g_{ij}- m(\max_{i \in I}f_{i}).


\end{equation}





Далее покажем, что функция $H(S), S \subseteq I$, является субмодулярной и, 


следовательно, можно воспользоваться жадным алгоритмом A, чтобы


получить приближенное решение задачи $(\ref{l21})$--$(\ref{l23})$


с гарантированной оценкой


точности, воспользовавшись теоремой \ref{t1}. 


B статье \cite{NWF78} было показано, что целевая функция 


задачи размещения с ограниченными объемами производств является 


субмодулярной. Здесь докажем более общий результат.





Рассмотрим задачу


\begin{gather*}


%\label{l31}


\max_{S \subseteq I} h(S), \\


%


%\label{l310}


h(S)= \max_{(x_{ij})}


\sum_{i \in S} \sum_{j \in J} (g_{ij} +F)x_{ij}, \quad


%\label{l32}


\sum_{i \in S}x_{ij} \le b_j,\ \; j \in J, \quad


%\label{l320}


\sum _{j \in J} x_{ij} \le d_i,\ \; i \in S, \\


%


%\label{l33}


0\le x_{ij}\le a_{ij},\ \; x_{ij}\; \text{ целое, \;}\ \; i\in S,\ \; j\in J.


\end{gather*}





\begin{lem}


\label{c5}


Функция $h(S)$, $S \subseteq I$, является субмодулярной.


\end{lem}





\begin{pf}


Представленная задача является задачей транспортного типа.


Рассмотрим следующую задачу с той же целевой функцией и теми же требованиями


на объемы поставок для каждого $j \in J$. 


Вместо множества $I$ рассмотрим множество $D$,


$|D|= \sum\limits_{i \in I} d_{i}$, 


пунктов отправления продукта с единичной мощностью каждый. Множество $D$ 


состоит из объединения непересекающихся множеств


$D_{i}$, $|D_{i}|= d_{i}$, $i \in I$, таких, что


источник из множества $D_{i}$ может поставлять продукт 


в любой пункт потребления, но при этом


не более $a_{ij}$ единиц из всех источников


$D_i$ в пункт потребления $j$. 





Целевая функция такой задачи является субмодулярной. 


Это следует из  эквивалентного определения субмодулярной 


функции \cite{NWF78}:


функция является субмодулярной, если 


$\rho_k(S) \ge \rho_k(T)$ 


$\forall S \subseteq T \subseteq D$, 


$\forall k \in D \setminus T$. 


Действительно, пусть $k \in D \setminus T$, $k \in D_{i}$. 


Если $h(S \cup \{k \}) - h(S) = g_{ij}$,


то $h(T \cup \{k \}) - h(T) \le g_{ij}$;


если же $h(S \cup \{k \}) - h(S) = 0$,


то $h(T \cup \{k \}) - h(T) = 0$, т.\,к. $S \subseteq T$.


Очевидно, целевая функция $h(T)$, $T \subseteq D$, является неубывающей. 





Пусть $h(S) = h( \cup_{i \in S}D_{i})$, $S \subseteq I$.


Тогда согласно лемме \ref{c4} функция $h(S)$, $S \subseteq I$,


является субмодулярной.


\end{pf}





Если положить $d_i=-f_i$, $i \in I$, то согласно леммам \ref{c2} и \ref{c3}


имеем следующее утверждение.





\begin{lem}


\label{c6}


Функция $H(S)$, $S \subseteq I$, является субмодулярной.


\end{lem}





Теперь можно сформулировать основной результат работы.





\begin{thm} 


Для решения задачи размещения на максимум с ограниченными объемами 


производств и поставок


$(\ref{l1})$--$(\ref{l3})$ существует полиномиальный алгоритм,


который находит приближенное решение $Z_G$


со следующей оценкой точности со сдвигом:


\begin{equation}


\frac{Z^* -Z_G}{Z^*- Z_R} < \frac{1}{e},


\end{equation}


где 


$Z^*$ --- точное решение, 


$Z_R$ имеет вид $(12)$, $e$ --- основание натурального логарифма.


\end{thm}





\begin{pf}


Применяем жадный алгоритм A к задаче $(\ref{l21})$--$(\ref{l23})$.


Можем воспользоваться теоремой $\ref{t1}$, т.\,к. по лемме \ref{c6}


целевая функция этой задачи субмодулярна и 


$H(S) \in C(\Theta)$, где $\Theta = \max\limits_{i \in I}f_i$.


Если в ходе этого алгоритма выбирается $k=m$ предприятий, то


получается (13).


Если выбрано $k<m$ предприятий, то относительная оценка со сдвигом


остается справедливой согласно цепочке неравенств


$$


\frac{Z^* -Z_G}{Z^*-Z(\emptyset)+m\Theta} \le 


\frac{Z^* -Z_G}{Z^*-Z(\emptyset)+k\Theta} 


\le (1 - 1/k)^k < 1/e.


$$


При этом точное решение $Z^*$ и допустимое решение $Z_G$ задачи


$(\ref{l21})$--$(\ref{l23})$ являются соответственно точным и допустимым 


решениями задачи $(\ref{l1})$--$(\ref{l3})$ согласно леммам \ref{l4} и


\ref{l41}.


Таким образом, найдено приближенное решение задачи


$(\ref{l1})$--$(\ref{l3})$ с гарантированной оценкой точности.





Заметим, что при вычислении $\rho_i(S)$ 


в ходе алгоритма A решается


транспортная задача на максимум с дополнительными ограничениями 


на объемы поставок.


Можно свести ее


к задаче о потоке минимальной стоимости следующим образом. 





Рассмотрим сеть с множеством вершин


$\{s \} \cup I \cup I^* \cup J \cup \{ t \}$, где $|I^*|= |I|$, и 


$s$, $t$ --- соответственно источник и сток сети.


Дуги ориентированы от $s$ к $i$ и от $j$ к $t$, 


$i \in I$, $j \in J$.


Каждый пункт $i \in I$ соединен выходящей дугой со вспомогательным


пунктом $i^* \in I^*$, из которого выходят дуги в каждый пункт


$j \in J$.


Пропускную способность дуги $(i, i^*)$ полагаем 


равной $d_i$, а дуги $( i^*, j)$ --- равной $a_{ij}$, 


остальные пропускные способности не ограничены.


Цену транспортировки единицы продукта из $i^*$ в $j$ полагаем равной 


$(-g_{ij})$, транспортировка по остальным дугам равна нулю.


Требуется найти поток величины $\sum\limits_{j \in J} b_j$ 


минимальной стоимости в данной сети.





Так как для задачи о потоке минимальной стоимости существует


полиномиальный алгоритм \cite{O93}, то


получаем полиномиальность построенного


жадного алгоритма A для задачи $(\ref{l1})$--$(\ref{l3})$.


\end{pf}





Выражаю признательность моему научному руководителю Гимади Э.Х. за 


помощь при написании этой статьи.





\begin{thebibliography}{99}





\bibitem{V99}


Вознюк И.П. 


{\em Задача размещения на сети с ограниченными пропускными способностями


коммуникаций} //


Дискретн. анализ и исследов. операций. Сер.\,2. -- 1999. -- 


Т.\,6. -- \No\,1. -- С.\,3--11.


\bibitem {GD82}


Гэри М., Джонсон Д. 


{\em Вычислительные машины и труднорешаемые задачи}. -- 


М.: Мир, 1982. -- 416 с.


\bibitem {H82}


Hochbaum D.S. 


{\em Heuristics for the fixed cost median problem} //


Math. Program. -- 1982. -- V.\,22. -- \No\,2. -- P.\,148--162.


\bibitem {STA97}


Shmoys D.B., Tardos E., Aardal K. 


{\em Approximation algorithms for facility location problems} //


In Proceedings of the 29th ACM Symposium on Theory of Computing. -- 


1997. -- P.\,265--274.


\bibitem {GK98}


Guha S., Khuller S.


{\em Greedy strikes back$:$ improved facility location algorithms} //


In Proceedings of the 9th Annual ACM-SIAM Symposium on Discrete 


Algorithms. -- 1998. -- P.\,649--657.


\bibitem {C98}


Chudak F.


{\em Improved approximation algorithms for uncapacitated


facility location} //


In Proceedings of the 6th IPCO Conference. -- 


1998. -- P.\,180--194.


\bibitem {KPR98}


Korupolu M., Plaxton C., Rajaraman R.


{\em Analysis of a local search heuristic for facility location problems} //


In Proceedings of the 9th Annual ACM-SIAM Symposium on Discrete 


Algorithms. -- 1998. -- P.\,1--10.


\bibitem {CW99}


Chudak F., Williamson D.


{\em Improved approximation algorithms for capacitated facility


location problems} //


In Proceedings of the 7th IPCO Conference. -- 


1999. -- P.\,99--113.


\bibitem {CFN77}


Cornuejols G., Fisher M.L., Nemhauser G.L.


{\em Location of bank accounts to optimize float\/$:$ an analytic study of


exact and approximate algorithms} //


Manag. Sci. -- 1977. -- V.\,23. -- \No\,8. -- P.\,789--810.


\bibitem {AS99}


Ageev A.A., Sviridenko M.I. 


{\em An $0.828$-approximation algorithm for the uncapacitated


facility location problem} //


Discrete Appl. Math. -- 1999. -- V.\,93. -- \No\,2--3. -- P.\,149--156.


\bibitem {NWF78}


Nemhauser G.L., Wolsey L.A., Fisher M.L. 


{\em An analysis of approximations for maximizing submodular set


functions} -- I //


Math. Program. -- 1978. -- V.\,14. -- \No\,3. -- P.265--294.


\bibitem {O93}


Orlin J.B.


{\em A faster strongly polynomial minimum cost flow algorithm} //


Operations research. -- 1993. -- V.\,41. -- \No\,2. -- P.\,338--350. 





\end{thebibliography}





\vskip 2cm





\post{\it Институт математики}{\it Поступила}


\post{\it Сибирского отделения}{05.07.1999}


\post{\it Российской Академии наук}{}





\label{end}


\end{document}


